\exercisegroup

\begin{enumerate}
\def\labelenumi{\arabic{enumi})}
\tightlist
\item
  \begin{enumerate}
  \def\labelenumii{\alph{enumii})}
  \tightlist
  \item
    Let E be a vector space. We say that a pointed convex cone C (of vertex 0) in E is maximal if C is a maximal element of the set of convex pointed cones of vertex 0 and \(\neq E\) , ordered by inclusion. Show that a pointed convex cone C is maximal if, and only if, it is a closed half-space defined by a hyperplane which passes through 0. To establish this result, prove successively the following properties of a maximal pointed convex cone C;
  \end{enumerate}
\end{enumerate}

α) We have C ∪ (− C) = E (argue by obtaining a contradiction).

\(\beta)\) If \(z\) is a non-internal (II, p.~26) point of C then \(-z \in C\) (same method). Deduce that C contains internal points.

\(\gamma)\) The largest vector subspace \(H = C \cap (-C)\) contained in C is a hyperplane. (Passing to the quotient space F = E/H, this reduces to demonstrating, using \(\beta)\) that if all the points of C other than the vertex are internal, then E is necessarily of dimension 1.)

\begin{enumerate}
\def\labelenumi{\alph{enumi})}
\setcounter{enumi}{1}
\tightlist
\item
  Give an example of a maximal non-pointed convex cone (in the set of non-pointed convex cones of vertex 0) which has no internal point (cf.~II, p.~65, exerc. 5).
\end{enumerate}

\begin{enumerate}
\def\labelenumi{\arabic{enumi})}
\setcounter{enumi}{1}
\item
  Let N be a hyperplane in a vector subspace M of a vector space E and let A be a convex set in E, such that all the points of \(A \cap M\) are on the same side of N and which also possesses the following property; for any \(y \neq 0\) in E, there exists \(x \in A \cap M\) such that \(x + \lambda y \in A\) for all \(\lambda\) such that \(|\lambda|\) is sufficiently small. Show that there exists then, a hyperplane H of E such that all the points of A are on the same side of H and such that \(H \cap M = N\) . (Reduce to the case \(N = \{0\}\) ; if \(a \neq 0\) belongs to \(A \cap M\) , consider the set U of pointed convex cones with vertex 0 containing A and not containing -a; show that there exists a maximal element C of U and that C is a maximal pointed convex cone (exerc. 1)). Deduce a new proof of the Hahn-Banach theorem.
\item
  Let A be a convex set in a topological vector space E and \(x_{0}\) be a point of E. Then, there exists a closed hyperplane H, containing \(x_{0}\) , and such that all the points of A lie on the same side of H if and only if there exists a non-pointed convex cone C with vertex \(x_{0}\) , which contains at least one interior point and does not meet A. (For an example of a convex set \(A \neq E\) which is not contained in any half-space defined by a hyperplane, see II, p.~65, exerc. 5.)
\item
  Let E be a normed space and A a complete convex set for the uniform structure induced by that of E.
\end{enumerate}

\begin{enumerate}
\def\labelenumi{\alph{enumi})}
\item
  Let \(x'\) be a continuous linear form on \(E\) that is bounded in \(A\) . Consider a number \(k > 0\) and the closed convex cone \(P\) in \(E\) , with vertex 0 and formed by the \(x \in E\) such that \(\| x \| \leqslant k < \langle x, x' \rangle\) ; it is pointed and proper. Show that for the order on \(E\) for which \(P\) is the set of elements \(\geqslant 0\) , the set \(A\) is inductive (use the fact that the restriction of \(x'\) to \(A\) is increasing and bounded).
\item
  Deduce from a) that the set of points of the frontier F of A which belong to a support hyperplane of A is dense in F (Bishop-Phelps th.). (For each point \(z \in F\) , consider a point \(y \in \complement A\) arbitrarily close to z, and separate y strictly from A by a closed hyperplane of equation \(\langle x, x' \rangle = \alpha\) , with \(\|x' \| = 1\) and use a) with k > 1, also exerc. 3 above.)
\end{enumerate}

\begin{enumerate}
\def\labelenumi{\arabic{enumi})}
\setcounter{enumi}{4}
\tightlist
\item
  Let \(\mathbf{A}\) be a closed convex set in \(\mathbf{R}^n\) and \(x_0\) be a point of \(\complement \mathbf{A}\) ; denote the euclidean distance in \(\mathbf{R}^n\) by \(d\) .
\end{enumerate}

\begin{enumerate}
\def\labelenumi{\alph{enumi})}
\item
  Show, without using th. 1 of II, p.~36, that there exists one and only one point \(x \in \mathbf{A}\) such that \(d(x_0, x) = d(x_0, \mathbf{A})\) , and that the hyperplane orthogonal to the line joining \(x_0\) and \(x\) , and passing through \(x\) is a support hyperplane of \(\mathbf{A}\) .
\item
  Deduce from a) a new proof of th. 1 of II, p.~36 when the space E is finite dimensional. (Reduce to the case when M is a frontier point \(x_{0}\) of A; note that the lower bound of the distance of \(x_{0}\) from support hyperplanes of A, is zero, and use the compactness of \(S_{n-1}\) .)
\end{enumerate}

\begin{enumerate}
\def\labelenumi{\arabic{enumi})}
\setcounter{enumi}{5}
\item
  Let A be a closed set in \(R^{n}\) with the following properties; for every \(x \in R^{n}\) , there exists one and only one point \(y \in A\) such that \(d(x, y) = d(x, A)\) , where d is the euclidean distance. Show that A is convex. (Argue by reductio ad absurdum, considering a closed segment with end points a, b in A containing a point \(c \in \complement A\) ; there is a closed ball B of centre c contained in \(\complement A\); consider the set B of closed balls S which contain B and whose interiors do not meet A; show that the radii of these balls is bounded above, and deduce that there exists one of these balls \(S_{0}\) whose radius \(\rho\) is the largest possible. Then get a contradiction by proving that \(S_{0}\) can only meet A in a single point, and that this implies the existence in B of a ball of radius \(> \rho\) .)
\item
  In a Hausdorff locally convex space \(\mathbf{E}\) , let \(\mathbf{A}\) be a complete, convex set and let \(\mathbf{B}\) be a precompact closed convex set such that \(\mathbf{A} \cap \mathbf{B} = \varnothing\) . Show that there exists a closed hyperplane separating A from B (argue in the completion \(\hat{\mathbf{E}}\) ). Consider the case when A is finite dimensional.
\item
  In a Hausdorff locally convex space, let A and B be two closed convex sets, without common points, such that \(C_{A} \cap C_{B} = \{0\}\) (II, p.~67, exerc. 14), and such that B is locally compact. Show that there exists a closed hyperplane separating A from B (cf.~II, p.~67, exerc. 16). Similarly, if A and B are two closed convex cones of vertex 0, such that \(A \cap B = \{0\}\) and B is locally compact, then there exists a closed hyperplane passing through 0 and separating A from B (use lemma 1 of II, p.~39).
\item
  Deduce from exerc. 8 that if V is a finite dimensional vector subspace of E and C is a closed convex cone of vertex 0 in E such that \(C \cap V = \{0\}\) , then there exists a support hyperplane of C that contains V (use lemma 1 of II, p.~39).
\item
  In the normed space \(E = l^{1}(\mathbf{N})\) of summable sequences of real numbers \(x = (\xi_{n})_{n \in \mathbf{N}}\) , let D be the line defined by the relations \(\xi_{n} = 0\) for \(n \geqslant 1\) . Show that there exist two increasing sequences \((\alpha_{n})\) , \((\beta_{n})\) of real numbers > 0 such that the convex set A defined by the inequalities \(\xi_{0} \geqslant |\alpha_{n}\xi_{n} - \beta_{n}|\) for \(n \geqslant 1\) is closed, non-bounded does not meet D and that there is no closed hyperplane separating A from D (choose \(\alpha_{n}\) and \(\beta_{n}\) so that A - D is everywhere dense).
\end{enumerate}

* 11) a) Let E be a Hilbert space and F an everywhere dense subspace of the dual E' of E that is distinct from E'; the unit ball B of E is compact for the weak topology σ(E, F), and there exists a point a of the unit sphere through which passes no closed (in σ(E, F)) support hyperplane of B.

\begin{enumerate}
\def\labelenumi{\alph{enumi})}
\setcounter{enumi}{1}
\item
  Give E the topology \(\sigma(\mathrm{E},\mathrm{F})\) and consider, in the product space \(\mathbf{G} = \mathbf{E} \times \mathbf{R}\) the set A of pairs \((x,\zeta)\) such that \(\|x\| < 1, \zeta \geqslant \|x\| / (1 - \|x\|)\) . Show that A is closed and locally compact, but that if D is the line with equation \(x = a\) in G then \(\mathrm{D} \cap \mathrm{A} = \emptyset\) and there does not exist any closed hyperplane in G separating A from D.
\item
  Show that, when we give E the topology \(\sigma(E, F)\) , there exists a continuous affine real valued function in the subspace B of E, that is not the restriction to B of a continuous affine function in E. *
\end{enumerate}

\begin{enumerate}
\def\labelenumi{\arabic{enumi})}
\setcounter{enumi}{11}
\item
  Consider in \(\mathbf{R}^3\) , the closed convex cone \(C\) defined by the relations \(\xi_1 \geqslant 0\) , \(\xi_2 \geqslant 0\) , \(\xi_3^2 \leqslant \xi_1\xi_2\) . Show that the line \(D\) of equations \(\xi_1 = 0\) , \(\xi_3 = 1\) does not meet \(C\) , but that there is no plane through the origin \(0\) containing \(D\) and not meeting \(C - \{0\}\) .
\item
  Let A and B be two closed convex sets in the space \(R^{n}\) , such that if V and W are affine linear varieties generated by A and B respectively, then no point of \(A \cap B\) is both interior to A relative to V and interior to B relative to W. Show that there exists a hyperplane separating A from B. (By taking quotients, reduce it to the case where either one of the varieties V, W is contained in the other or V and W are complementary vector subspaces in E.)
\item
  Let A be a parabolic closed convex set (II, p.~67, exerc. 17) not containing a line. Show that if B is a closed convex set not meeting A then there exists a hyperplane in \(\mathbf{R}^n\) that separates A strictly from B (if \(d\) is the Euclidean distance prove that \(d(\mathrm{A},\mathrm{B}) > 0\) ) (cf.~exerc. 12.)
\item
  Let S, T two finite sets in \(R^{n}\) with no common points, and such that \(\operatorname{Card}(S \cup T) \geqslant n + 2\) . In order that there exists a hyperplane separating S strictly from T, it is necessary and sufficient that for every finite set \(F \subset S \cup T\) of \(n + 2\) points, there exists a hyperplane separating \(F \cap S\) strictly from \(F \cap T\) (use Helly's th. (II, p.~68, exerc. 21)). Show that in this statement we cannot replace the number \(n + 2\) by \(n + 1\) , and that the statement does not extend to the case where S and T are infinite.
\item
  Let A be a compact set with interior points in \(R^{n}\) . Show that if each frontier point of A lies on at least one support hyperplane of A, then A is convex. (Obtain a contradiction, showing that if x and y are two points of A such that the segment with end points x, y is not contained in A, and if z is an interior point of A not situated on this segment, then there exists a frontier point of A, distinct from x and y in the triangle with vertices x, y, z.)
\item
  In \(R^{n}\) , let A be a symmetric convex set of which 0 is an interior point and of which the frontier does not contain any genuine segment. Let H be a homogeneous hyperplane and D a line complementary to H. Show that there exists a point \(a \in H \cap A\) such that at a there is a hyperplane of support to A that is parallel to D.
\item
  In a topological vector space \(\mathbf{E}\) , let \(\mathrm{A}_i\) ( \(1 \leqslant i \leqslant n\) ) be \(n\) open non-empty convex sets. a) Show that if the union of the \(\mathrm{A}_i\) is distinct from \(\mathbf{E}\) , then every point \(x \in \mathbf{E}\) not belonging to any of the \(\mathrm{A}_i\) , belongs to a closed linear variety of codimension \(n\) , that contains \(x\) and does not meet any of the \(\mathrm{A}_i\) (argue by induction on \(n\) ).
\end{enumerate}

\begin{enumerate}
\def\labelenumi{\alph{enumi})}
\setcounter{enumi}{1}
\tightlist
\item
  If the intersection of the \(\mathbf{A}_i\) is empty, show that there exists, in E, a closed linear variety of codimension \(n - 1\) that does not meet any of the \(\mathbf{A}_i\) (same method).
\end{enumerate}

\begin{enumerate}
\def\labelenumi{\arabic{enumi})}
\setcounter{enumi}{18}
\tightlist
\item
  Let C, C' be two closed convex sets in a Hausdorff topological vector space E that are strictly separated by a closed hyperplane H. Let H' be a closed hyperplane of support to both C and C' such that C and C' lie on the same side of H'. Show that H' is the only hyperplane with these properties which contains H ∩ H' and that H ∩ H' is a support hyperplane of the trace P on H of the convex envelope of C ∪ C'. Conversely, if C and C' are compact, then for every support hyperplane D of P in H, there exists a hyperplane H' that supports both C and C', which contains D and such that C and C' lie on the same side of H'.
\end{enumerate}

¶ 20) Let A, B be two disjoint closed convex sets in a Hausdorff locally convex space E and let H be a closed hyperplane separating A from B; suppose that \(A \cap H \neq \varnothing\) and that the intersection of \(A \cap H\) and of every line is compact. Show that, if A or B is locally compact, then there exists a neighbourhood V of 0 in E such that \((A + V) \cap B\) is empty. (Consider two cases according to whether A or B is locally compact; in the first case, note that there exists a hyperplane \(H'\) parallel to H such that, if S is the set of points between H and \(H'\) , then \(A \cap S\) is compact. In the second case, suppose for example that \(0 \in B \cap H\) ; for every neighbourhood V of 0 in E, consider the set \((A + V) \cap B\) and consider successively the case where this set is relatively compact for at least one V or the case when this is not so, as in exerc. 16 of II, p.~67.)

¶ 21) a) In \(\mathbf{R}^n\) let \(a_i (1 \leqslant i \leqslant n + 1)\) be \(n + 1\) points that are affinely independent. Denote the convex envelope of the \(a_i\) by \(\mathbf{S}\) and the convex envelope of the \(a_i\) with \(i \neq k\) by \(\mathrm{F}_k\) for \(1 \leqslant k \leqslant n + 1\) . For each \(k\) let \(\mathrm{C}_k\) be a compact convex set containing \(\mathrm{F}_k\) , and suppose that

S is contained in the union of the \(C_k\) ; show then that \(\bigcap_{k=1}^{n+1} C_k \neq \emptyset\) . (Argue by reductio ad

absurdum and induction on \(n\) , considering the intersection \(\mathbf{C}_{n+1}'\) of the \(\mathbf{C}_i\) with indices \(i \leqslant n\) and supposing that \(\mathbf{C}_{n+1} \cap \mathbf{C}_{n+1}' = \emptyset\) , which would allow the strict separation of the two convex sets by a hyperplane.)

\begin{enumerate}
\def\labelenumi{\alph{enumi})}
\setcounter{enumi}{1}
\tightlist
\item
  Let X be a compact convex set in a Hausdorff topological vector space E, and \((\mathrm{C}_{\lambda})_{\lambda\in\mathrm{L}}\) a family of compact convex sets contained in X, such that for every set H ⊂ L having n (resp. m) elements, the intersection (resp. the union) of the \(C_{\lambda}\) with indices \(\lambda\in H\) is not empty (resp. is equal to X). Show that if \(m\leqslant n+1\) the intersection \(\bigcap_{\lambda\in L}C_{\lambda}\) is not empty. (This is effectively
\end{enumerate}

proving that for any finite set H of \(p \geqslant m\) indices of L, we have \(\bigcap_{\lambda \in H} C_{\lambda} \neq \varnothing\) . Argue by induc-

tion on \(p\) assuming that the result has been proved for \(p - 1\) indices. Argue then by reductio ad absurdum, considering for each index \(i \in \mathbf{H}\) a point \(a_i \in \bigcap_{\lambda \in \mathbf{H} - \{i\}} \mathrm{C}_\lambda\) , and showing by the

aid of Helly's th. (II, p.~68, exerc. 21) that the \(a_{i}\) generate a linear variety of dimension \(p - 1\) , then finally apply \(a\) .)

¶ 22) In \(\mathbf{R}^n\) let \((\mathbf{C}_i)_{1 \leqslant i \leqslant m}\) be a finite family of closed convex cones with vertex 0, such that the sum of any \(n\) of them is distinct from \(\mathbf{R}^n\) . Show that there exists a hyperplane H, passing through 0, such that, for any index \(i\) no pair of points of \(\mathbf{C}_i\) are strictly separated by H. (Distinguish two cases:

\(\alpha)\) Either there exists a number \(r < n\) and \(r\) indices, say, 1, 2, \ldots, \(r\) such that \(C_i\) for \(i \leqslant r\) generates a cone which contains a vector subspace V of dimension \(\geqslant r\) . Argue by induction on \(n\) , projecting on the orthogonal to V.

β) Or, for all r < n, any r of the \(C_{i}\) generate a cone C such that the maximal vector subspace \(C \cap (-C)\) contained in C is of dimension < r. Consider then a set of the cones \(C_{i}\) , maximal with respect to being contained in a half-space; there are at least n cones in a maximal set. Let \(\Gamma\) be the cone generated by the union of the cones belonging to this maximal set. If \(C_{j}\) is a cone which does not belong to the maximal set considered, show that \(C_{i} \subset -\Gamma\) . For this, argue by reductio ad absurdum, showing that in the contrary case there exists a frontier point of \(-\Gamma\) (relative to the vector subspace generated by \(\Gamma\) ) that is interior to \(C_{j}\) (relative to the vector subspace generated by \(C_{j}\) ). Write such a point as the sum of the least number s of vectors, of which each belongs to a cone \(-C_{i}\) , among those \(C_{i}\) used in defining \(\Gamma\) ; then \(s \leqslant n - 1\) . Prove finally that these cones and \(C_{j}\) generate a convex cone containing an \(s + 1\) dimensional vector subspace, contradicting the hypothesis; for this use exerc. 4 of II, p.~65.)

\begin{enumerate}
\def\labelenumi{\arabic{enumi})}
\setcounter{enumi}{22}
\item
  Let \(\mathbf{E}\) be a topological vector space, and let \(\mathcal{T}\) be the locally convex topology on \(\mathbf{E}\) that is the finest of all those that are coarser than the given topology \(\mathcal{T}_0\) on \(\mathbf{E}\) . If \(\mathbf{F}\) is a locally convex space, then the continuous linear mappings of \(\mathbf{E}\) in \(\mathbf{F}\) are the same for \(\mathcal{T}_0\) as for \(\mathcal{T}\) . There exists a continuous linear form on \(\mathbf{E}\) that is distinct from the null form if, and only if, there exists a neighbourhood of 0 for \(\mathcal{T}_0\) whose convex envelope is not everywhere dense (for \(\mathcal{T}_0\) ) (cf.~I, p.~25, exerc. 4).
\item
  Let \(\mathbf{E}\) be an infinite dimensional, metrisable, locally convex space.
\end{enumerate}

\begin{enumerate}
\def\labelenumi{\alph{enumi})}
\item
  Show that there exists a sequence \((a_{n})\) of points of E tending to 0 and a decreasing sequence \((\mathbf{L}_n)\) , of closed vector subspaces of E, such that \(\mathbf{L}_n\) is of codimension \(n\) in E and that, for all \(n\) , the point \(a_{n}\) belongs to \(\mathbf{L}_n - \mathbf{L}_{n + 1}\) .
\item
  Suppose further that E is complete. Show that we can then find sequences \((a_{n})\) and \((\mathrm{L}_{n})\) verifying the conditions a), and such that in addition, for every bounded sequence of real number \((\lambda_{n})\) , the series, whose general term is \(\lambda_{n}a_{n}\) , is commutatively convergent in E, and that the linear mapping \((\xi_{n})\mapsto\sum_{n}\xi_{n}a_{n}\) of the Banach space \(\mathcal{B}(\mathbf{N})\) in E is injective and continuous.
\item
  Deduce from b) that when E is an infinite dimensional Fréchet space then every basis of E on R has cardinal at least equal to \(2^{\mathrm{Card}(\mathbf{N})}\) (cf.~I, p.~22, exerc. 5).
\end{enumerate}

If there exists an enumerable set that is dense in \(\mathbf{E}\) , then every basis of \(\mathbf{E}\) has the cardinal of the continuum.

\begin{enumerate}
\def\labelenumi{\arabic{enumi})}
\setcounter{enumi}{24}
\item
  Let E be an infinite dimensional Fréchet space of enumerable type (therefore having an enumerable everywhere dense subset) (cf.~I, p.~25, exerc. 1). Show that there exists an everywhere dense hyperplane H of E which meets every closed, infinite dimensional linear variety of E. (Use the existence of a basis, having the cardinal of the continuum, in each of the direction subspaces of these varieties (exerc. 24, c)) and the fact that the set of closed, infinite dimensional, linear varieties of E also has the cardinal of the continuum (GT, IX, § 5, exerc. 17); then apply a method of construction of a linear form on E following from S, III, § 6, exerc. 24.) The hyperplane H does not contain any infinite dimensional, closed, vector subspace.
\item
  Let \(\mathbf{E}\) be an infinite dimensional Fréchet space of enumerable type.
\end{enumerate}

\begin{enumerate}
\def\labelenumi{\alph{enumi})}
\item
  Show that there exists a sequence \((a_{n})\) of linearly independent elements of E such that each sequence \((a_{2n})\) and \((a_{2n + 1})\) is total (use exerc. 24, c)).
\item
  Let \(\mathbf{F}\) be the vector subspace of \(\mathbf{E}\) generated by the \(a_{2n + 1}\) ( \(n \in \mathbf{N}\) ). For every \(n > 0\) let \(\mathbf{M}_n\) be the subspace generated by the \(a_{2k}\) with \(k \leqslant n\) . For each \(n\) , let \(\phi_n\) be the restriction to \(\mathbf{F}\) of the canonical homomorphism of \(\mathbf{E}\) on \(\mathrm{E / M}_n\) , and let \(\mathcal{T}_n\) be the topology on \(\mathbf{F}\) which is the inverse image under \(\phi_n\) of the quotient topology, on \(\mathrm{E / M}_n\) . Show that each of the topologies \(\mathcal{T}_n\) on \(\mathbf{F}\) is a Hausdorff locally convex topology, but that the lower bound of the \(\mathcal{T}_n\) in the set of locally convex topologies on \(\mathbf{F}\) is the coarsest topology on \(\mathbf{F}\) .
\end{enumerate}

* c) Take E to be a Hilbert space; show that we can choose the sequence \((a_{n})\) so that if G is the closed vector subspace generated by the \(a_{4n + 1}\) , then G has infinite codimension and so that the images of the \(a_{2n}\) and the \(a_{4n + 3}\) in E/G are still linearly independent. Write \(\mathbf{G}_n\) for the subspace of E that is the sum of G and of the subspace generated by the \(a_{4k + 3}\) with \(k\leqslant n\) , and give to \(G_{n}\) the topology which is the inverse image under the canonical mapping restricted to \(G_{n}\) , of the quotient topology on \(E/M_{n}\) . Show that the sequence \((\mathbf{G}_{n})\) is an inductive system of topological vector spaces such that \(G_{n}\) is closed in \(G_{n+1}\) for the topology of \(G_{n+1}\) , but that \(G_{n}\) is not closed in the inductive limit space of this sequence. *

\begin{enumerate}
\def\labelenumi{\arabic{enumi})}
\setcounter{enumi}{26}
\tightlist
\item
  Let E, F be two Hausdorff topological vector spaces, and X (resp. Y) a compact convex set in E (resp. F). Let f be a real valued function defined in \(X \times Y\) with the following properties:
\end{enumerate}

\begin{enumerate}
\def\labelenumi{(\roman{enumi})}
\item
  For all \(x \in \mathbf{X}\) , the mapping \(y \mapsto f(x, y)\) is lower semi-continuous in \(\mathbf{Y}\) , and for all \(c \in \mathbf{R}\) , the set of the \(y \in \mathbf{Y}\) such that \(f(x, y) \leqslant c\) is convex.
\item
  For all \(y \in \mathbf{Y}\) , the mapping \(x \mapsto f(x, y)\) is upper semi-continuous in \(\mathbf{X}\) , and for all \(c \in \mathbf{R}\) , the set of the \(x \in \mathbf{X}\) such that \(f(x, y) \geqslant c\) is convex.
\end{enumerate}

Show that, in these conditions, we have

\[
\sup _ {x \in X} (\inf _ {y \in Y} f (x, y)) = \inf _ {y \in Y} (\sup _ {x \in X} f (x, y)).
\]

(Argue by reductio ad absurdum, supposing that there exists a number \(c\) such that

\[
\sup _ {x \in \mathbf {X}} (\inf _ {y \in \mathbf {Y}} f (x, y)) <   c <   \inf _ {y \in \mathbf {Y}} (\sup _ {x \in \mathbf {X}} f (x, y)).
\]

For all \(x \in X\) (resp. all \(y \in Y\) ) let \(A_{x}\) be the set of \(y \in Y\) such that \(f(x, y) > c\) (resp. \(B_{y}\) the set of the \(x \in X\) such that \(f(x, y) < c\) ), which is open in Y (resp. in X); the \(A_{x}\) (resp. the \(B_{y}\) ) form a covering of Y (resp. X) when x varies in X (resp. y varies in Y). Show that there exist two finite sets \(X_{0} \subset X\) , \(Y_{0} \subset Y\) such that: \(1^{\circ}\) for all y belonging to the convex envelope \(B_{0}\) of \(Y_{0}\) , there exists \(x \in X_{0}\) such that \(f(x, y) > c\) , and \(X_{0}\) is minimal for this property; \(2^{\circ}\) for all x belonging to the convex envelope \(A_{0}\) of \(X_{0}\) , there exists \(y \in Y_{0}\) such that \(f(x, y) < c\) , and \(Y_{0}\) is minimal for this property. Then for all \(y \in Y_{0}\) , let \(C_{y}\) be the set of \(x \in A_{0}\) such that \(f(x, y) \geqslant c\) ; using exerc. 21, a) of II, p.~79, show that the intersection of the \(C_{y}\) for \(y \in Y_{0}\) is not empty. Proceed in the same way in \(B_{0}\) and obtain a contradiction.)

¶ 28) Let X be a compact convex subset of E a Hausdorff locally convex space, and let f be an upper semi-continuous convex function in X. Show that the set L of continuous convex functions g in X such that \(g(x) > f(x)\) for all \(x \in X\) is decreasing directed and that its lower envelope is equal to f.~(Let u, v be elements of L. To construct an element of L which is less than u and v, use reasoning analogous to that of prop. 6, II, p.~40. Interpret the set \(K_{1}\) analogous to the set K in this argument as the set of points situated above the graph of a lower semi continuous function that is less than u and v and strictly larger than f at every point; apply prop. 5 of II, p.~39 and Dini's th. to this function. To show that the lower envelope of L is f, note that f is bounded above by a constant b; \((x, t)\) being a point of E × R situated above the graph of f, let \(K'\) be the convex envelope of \(\{(x, t)\} \cup (X' \times \{b\})\) , where \(X'\) is a convenient compact neighbourhood of x in X; argue with \(K'\) as above for \(K_{1}\) .)

\begin{enumerate}
\def\labelenumi{\arabic{enumi})}
\setcounter{enumi}{28}
\item
  Let X be a compact convex set in a Hausdorff locally convex space E. Let u be a lower semi-continuous convex function in X and v an upper semi-continuous concave function in X such that \(u(x) > v(x)\) for all \(x \in X\) . Then there exists an affine linear function f that is continuous in E and such that \(v(x) < f(x) < u(x)\) for all \(x \in X\) .
\item
  Let \(\mathbf{X}\) be a compact convex set of a Hausdorff locally convex space \(\mathbf{E}\) . Show that the set of lower semi-continuous convex functions in \(\mathbf{X}\) is a lattice.
\end{enumerate}
% semantic-fragments: legacy-input-seam
\relax{}
